
\begin{center}
	\Huge
	Talsystemer
\end{center}
\stepcounter{section}
Når vi repræsenterer heltal, så gør vi det med 10 symboler - vi regner i base 10. Dette har historiske grunde og er formentlig fordi vi har 10 fingre, men der er ingen grund til at vi kunne anvende andre talsystemer. Babylonerne regnede i 60-talssystemet, og dette er grunden til, at vi i dag har 60 minutter på en time og 360 grader i en cirkel. Vi skal nu se, hvordan vi kan skifte mellem talsystemer og regne i andre talsystemer. 

\begin{exa}
	Tallet $n = 274$ anvender symbolerne 2, 7 og 4. Den måde, vi forstår tallet er
	\begin{align*}
		n = 2\cdot 10^2 + 7\cdot 10^1 + 4\cdot 10^0.
	\end{align*}
	Det er altså tallet 10 her, der afslører, at vi arbejder i base 10. 
	For at notere, at et tal er i base 10, skrives der når man kan være i tvivl $274_{10}$.
	Vi kunne tilsvarende have tallet $m = 2432_{5}$. Dette forstås som
	\begin{align*}
		m = 2\cdot 5^3 + 4\cdot 5^2 + 3\cdot 5^1 + 2\cdot 5^0
	\end{align*}
	i 10-talssystemet. Udregner vi dette i titalssystemet, vil $m = 2432_{5}$ altså tilsvare
	\begin{align*}
		m &= 2\cdot 5^3 + 4\cdot 5^2 + 3\cdot 5^1 + 2\cdot 5^0 \\
		  &= 2\cdot 125 + 4\cdot 25 + 3\cdot 5 + 2\cdot 1 \\ 
		  &= 250 + 100 + 15 + 2 \\
		  &= 367_{10}.
	\end{align*}
	I femtalssystemet har vi kun symbolerne 0, 1, 2, 3 og 4.
\end{exa}

Vi kan opsummere dette i følgende sætning.
\begin{setn}
	Ethvert positivt heltal $n\in \mathbb{Z}$ kan repræsenteres entydigt i base $b\in \mathbb{Z}_{>1}$ som
	\begin{align*}
		n = a_k b^k + a_{k-1}b^{k-1} + \cdots + a_1 b + a_0,
	\end{align*}
	hvor $a_1, a_2, \hdots a_k \in \mathbb{Z}/b\mathbb{Z}$ og $a_k \neq 0$. 
\end{setn}

\begin{exa}
	De \textit{binære} tal har en særlig rolle i digital kommunikation, da et elektrisk signal (i al praktisk henseende)
	 kun kan være slukket og tændt. Vi anvender symbolerne 0 og 1 til at beskrive disse situationer henholdsvis. Lad os
	 betragte tallet $(10001011)_{2}$ i base 10. Vi får
	 \begin{align*}
	 	1\cdot 2^7 + 1\cdot 2^3 + 1\cdot 2 + 1 = 128 + 8 + 2 + 1 = 141_{10}
	 \end{align*}
\end{exa}

Det er også muligt at finde tal i titalssystemet (Decimaltal) i andre talsystemer. Dette er blot lidt mere besværligt. Vi betragter et eksempel.

\begin{exa}
	Vi ønsker at skrive 673 i base 6. Vi gør det ved at dividere med rest.
	\begin{align*}
		673 = 112 \cdot 6 + 1
	\end{align*}
	Vi deler nu kvotienten 112 med 6.
	\begin{align*}
		112 = 18\cdot 6 + 4.
	\end{align*}
	Dette gøres to gange til
	\begin{align*}
		18 &= 3 \cdot 6 + 0\\
		3 &= 0 \cdot 6 + 3
	\end{align*}
	Tallet er nu resterne, så 
	\begin{align*}
		673_{10} = 3041_{6}.
	\end{align*}
	Lad os tage endnu et eksempel. 
	Vi skriver 1000 i base 2.
	\begin{align*}
		1000 &= 500\cdot 2 + 0\\
		500 &= 250 \cdot 2 + 0\\
		250 &= 125 \cdot 2 + 0\\
		125 &= 62 \cdot 2 + 1\\
		62 &= 31\cdot 2 + 0\\
		31 &= 15\cdot 2 + 1\\
		15 &= 7\cdot 2 + 1\\
		7 &= 3\cdot 2 + 1\\
		3 &= 1\cdot 2 + 1\\
		1 &= 0\cdot 2 + 1
	\end{align*}
	Derfor er 
	\begin{align*}
		1000_{10} = 1111101000_{2}
	\end{align*}	
\end{exa}

For en god ordens skyld kan vi lige prøve at tælle til 10 i base 3. Dette kan ses på Figur \ref{fig:base3}
\begin{table}[H]
	\centering
	\rowcolors{1}{Melon!30}{SeaGreen!30}
	\begin{tabular}{c|ccccccccccc}
		Base 10&1&2&3&4&5&6&7&8&9&10 \\
		Base 3&1&2&10&11&12&20&21&22&100&101
	\end{tabular}
	\caption{Tæl til 10 i base 3}
	\label{fig:base3}
\end{table}

\section*{Addition og multiplikation}
\stepcounter{section}

Vi vil også gerne kunne lægge tal sammen og gange dem sammen i andre baser. Som før tager vi udgangspunkt i et eksempel.
\begin{exa}	
	Vi ønsker at lægge tallene $432_6$ og $1025_6$ sammen. Dette gøres som vi er vant til. Først lægges 2 og 5 sammen
	\begin{align*}
		2 + 5 = 7
	\end{align*}
	Dette bestemmes så $\fmod{6}$ og vi tager kvotienten med videre.
	\begin{align*}
		7 \fmod{6} = 1
	\end{align*}
	og kvotienten er 1, så vi får nu for næste ciffer én i mente
	\begin{align*}
		2 + 3 + 1 \fmod{6} = 0,
	\end{align*}
	og vi får igen én i mente. For tredje ciffer
	\begin{align*}
		4 + 0 + 1 \fmod{6} = 5,
	\end{align*}
	og 0 i mente. 
	Det 4. ciffer
	\begin{align*}
		1 + 0  \fmod{6} = 1
	\end{align*}
	Derfor er summen 
	\begin{align*}
		1501_6
	\end{align*}
\end{exa}
\begin{exa}
	Vi ønsker nu at gange $442_6$ med $12_6$. Fremgangsmåden er den samme som når vi ganger normalt, men som før
	 nu$\fmod{6}$ og vi får
	 \begin{align*}
	\begin{array}{ccccc}
		& &	4&4&2\\
		&\times & & 1&2\\
		\hline
		& & & & 4 \\
		& & 1&2 &0 \\
		&1 &2 &0 &0 \\
		& & & 2 & 0 \\
		& &4 &0 &0 \\ 
		& 4& 0& 0&0\\
		\hline
		1& 0& 1&4 &4
	\end{array}
	\end{align*}
	Der gælder altså derfor, at 
	\begin{align*}
		442_6 \cdot 12_6 = 10144_6
	\end{align*}
\end{exa}


\subsection*{Opgave 1}
Udfyld Tabel \ref{tab:tælle}.
\begin{table}[H]
	\centering
	\rowcolors{1}{SeaGreen!30}{Melon!30}
	\begin{tabular}{ccccccccccccccccc}
		Base 10 & 1 &2 &3&4&5&6&7&8&9&10&11&12 &13 &14 &15 &16 \\
		Base 4 &&&&&&&&&&&&&&&& \\
		Base 2 &&&&&&&&&&&&&&&& \\
		Base 8 &&&&&&&&&&&&&&&&
	\end{tabular}
	\caption{Tælle i andre baser}
	\label{tab:tælle}
\end{table}

\subsection*{Opgave 2}
Udregn følgende tal i decimaltal (base 10).
\begin{align*}
	&1) \ 55_6  & &2) \ 1011_2   & &3) \ 233_4  \\
	&4) \ 11_9  & &5) \ 210_3   & &6) \ 11111_2    \\ 
	&7) \ 241_5  & &8) \ 1111_3   & &9) \ 1010101_2   \\
\end{align*}

\subsection*{Opgave 3}
Omregn følgende tal fra base 10 til den angivne basis
\begin{align*}
	&1) \ 103 \textnormal{ i base 6}  & &2) \ 200 \textnormal{ i base 2}   & &3) \ 450 \textnormal{ i base 5}  \\
	&4) \ 256 \textnormal{ i base 2}  & &5) \ 621 \textnormal{ i base 8}   & &6) \ 172 \textnormal{ i base 9}   \\ 
	&7) \ 411 \textnormal{ i base 3}  & &8) \ 700 \textnormal{ i base 7}   & &9) \ 188 \textnormal{ i base 4}   \\
\end{align*}

\subsection*{Opgave 4}
Bestem følgende summer eller produkter
\begin{align*}
	&1) \ 112_4 + 211_4   &&2) \ 1011_2+1001_2    &&3) \  410_5\cdot 14_5  \\
	&4) \ 55_6\cdot 22_6   &&5) \  101_2\cdot 122_2   &&6) \ 77_8+504_8   \\
\end{align*}

\subsection*{Opgave 5}
Vi ønsker at bestemme summen $211_3 + 1002_3$.
\begin{enumerate}[label=\roman*)]
	\item Udregn summen $211_3 + 1002_3$.
	\item Omregn $211_3$ og $1002_3$ til titalsssystemet.
	\item Bestem summen af tallene
	\item Regn tilbage til base 3 og tjek dit svar.
\end{enumerate}

\subsection*{Opgave 6}
Der er ingen grund til, at vi ikke kan have flere symboler end 10. I det hexadecimale talsystem har man 16 symboler. Det er de sædvanlige fra titalssystemet tilføjet A, B, C, D, E, og F så $A_{16} = 11_{10}$, $B_{16} = 12_{10}$ osv.
\begin{enumerate}[label=\roman*)	]
	\item Opskriv tallene 15, 20, 66 og 102 i 16-talssystemet.
	\item Bestem summen $B12_{16}+2DE_{16}$. 
	\item Bestem produktet $34_{16}\cdot  FD_{16}$.
\end{enumerate}

\subsection*{Opgave 7}
Når vi i kryptologi skal bestemme koder, så skal vi kunne udregne udtryk som $3^{2114} \fmod{7}$. (I virkeligheden kan potenserne være meget større).
\begin{enumerate}[label=\roman*)]
	\item Opskriv en multiplikationstabel for $\mathbb{Z}/7\mathbb{Z}$. 
	\item Udregn $3^{2114} \fmod{7}$ ved gentagende at kvadrere tallet. 
\end{enumerate}

\subsection*{Opgave 8}
Udregn følgende
\begin{align*}
	&1) \ F21_{16} + 1011011_{2}   &&2) \ ABBA_{16} +  7712_{8}   &&3) \  1323_4 + 746_8  \\
\end{align*}